% (C) 2026 Ralf Stephan, in collaboration with Claude Code.  CC0 1.0.
\chapter{Covers of dynamical systems}\label{ch:covers}

So far every theorem has been about a graph. This chapter connects graphs to orbits. One
definition, the \emph{cover}, states what every ``the real orbit is a path'' lemma in the
corpus proves. Once it holds, the graph theory of Chapters~\ref{ch:core}--\ref{ch:splitting}
turns into statements about orbits with no further work: periodicity of carry words, escape
times, word counts, dimension bounds, and the obstruction principle. This chapter proves these
\emph{transfer theorems} once, for an abstract system. Chapter~\ref{ch:affine} supplies the
covers of the recurrence~\eqref{eq:carry}.

\section{Systems, runs and covers}

\begin{definition}\label{def:system}
A \emph{system}\index{system} over the alphabet \(A\) is a set \(S\) with maps \(f_a\colon S\to S\)
(\(a\in A\)) and an \emph{admissibility}\index{admissibility} relation \(\mathrm{adm}(x,a)\): ``the step \(a\) may be
taken at \(x\)''. A point \(x\) \emph{follows}\index{follow (a run)} the word \(a_1\cdots a_n\) if each step of the
\emph{run}\index{run} \(x_0=x\), \(x_k=f_{a_k}(x_{k-1})\) is admissible: \(\mathrm{adm}(x_{k-1},a_k)\) for
all \(k\). Write \(\mathrm{run}(x,w)=x_n\). A label sequence \(\ell\colon\N\to A\) is
\emph{admissible forever}\index{admissible forever} at \(x\) if \(x\) follows every prefix of \(\ell\).
\end{definition}

The maps \(f_a\) may ignore their label. In Chapter~\ref{ch:affine} there is a single map
\(z\mapsto(p/q)(z-\nu)+\nu\), and the label is the carry, which admissibility pins down. Avoiding
a divisor set is part of admissibility, so ``admissible forever'' means ``avoids \(\mathcal P\)
forever''.

\begin{definition}\label{def:cover}
A carry graph \(X\) on \(V\) \emph{covers}\index{cover} the system through \(\rho\colon V\to\mathcal P(S)\)
if every admissible step from \(\rho(v)\) is shadowed by an arc:
\[
x\in\rho(v),\ \mathrm{adm}(x,a)\ \Longrightarrow\ f_a(x)\in\rho(w)\text{ for some }
v\xrightarrow{a}w.
\]
\(X\) \emph{realises}\index{realise} its words if every word of \(X\) from \(v\) is followed by some point of
\(\rho(v)\), and \emph{realises forever}\index{realise!forever} if every infinite labelled path
\((s_n,\ell_n)\) of \(X\) is admissible forever at some point of \(\rho(s_0)\).
\end{definition}

A cover is a soundness statement and needs no converse. Upper bounds (periodicity, escape,
entropy, dimension) need only a cover. Lower bounds need realisation. A cover that also
realises its words is \emph{exact}: its words are the itineraries
(Theorem~\ref{thm:count}\,(b)).

\begin{proposition}[shadowing]\label{prop:shadow}\index{shadowing}
Let \(X\) cover the system through \(\rho\) and \(x\in\rho(v)\).
\begin{enumerate}[label=(\alph*)]
\item If \(x\) follows \(w\), then \(X\) reads \(w\) from \(v\) to some \(u\) with
  \(\mathrm{run}(x,w)\in\rho(u)\).
\item If \(\ell\) is admissible forever at \(x\), there is an infinite labelled path
  \((s_n,\ell_n)\) of \(X\) with \(s_0=v\) and \(x_n\in\rho(s_n)\) for all \(n\).
\end{enumerate}
\end{proposition}

\begin{proof}
(a) Induction on \(w\), one application of the cover per letter. (b) Choose \(s_{n+1}\) among
the arcs that the cover provides at step \(n\). The choice is recursive along the orbit, so no
compactness is needed and \(V\) may be infinite.
\end{proof}

\paragraph{Constructions.} Covers survive the operations of Chapter~\ref{ch:splitting}.
\begin{itemize}
\item A supergraph of a cover is a cover, and so is a cover of a system with fewer admissible
  steps.
\item The image of a cover under a homomorphism \(g\) is a cover through the fibre unions
  \(\bigcup_{g(v)=u}\rho(v)\).
\item If \(\pi\colon S'\to S\) maps admissible steps to admissible steps
  (\(\pi(f'_a x)=f_a(\pi x)\)), a cover of \(S\) covers \(S'\) through the preimages
  \(\pi^{-1}\rho(v)\). This is how the fractional parts enter in Chapter~\ref{ch:affine}.
\item Covers by \(X\) and \(Y\) give a cover by the label product \(X\otimes Y\) through
  \(\rho(v)\cap\sigma(w)\), and likewise for indexed products. This is the
  \emph{residue\,\(\times\)\,cell principle}: the residue and the cell of \(x_n\) are tracked
  separately and combined at no cost.
\item The subset construction covers through \(\rho(T)=\bigcup_{v\in T}\rho(v)\).
\item The line graph covers through \(\rho(u\xrightarrow{a}v)=\rho(v)\), and by iteration so does
  every higher block graph \cite[§2.3]{LM95}.
\end{itemize}

\section{The transfer theorems}

\begin{theorem}[periodicity]\label{thm:transfer-per}
Let \(X\) cover the system through \(\rho\), and let \(\ell\) be admissible forever at a point
of some \(\rho(v)\).
\begin{enumerate}[label=(\alph*)]
\item If \(X\) has the label form (Definition~\ref{def:allper}), \(\ell\) is eventually
  periodic.
\item In particular this holds if \(V\) is finite and \(X\) has zero entropy, or if \(\det X\)
  has zero entropy (Theorem~\ref{thm:determinization}).
\item Conversely, one admissible-forever \(\ell\) that is not eventually periodic rules out the
  label form of every cover.
\end{enumerate}
\end{theorem}

\begin{proof}
Proposition~\ref{prop:shadow}\,(b) makes \(\ell\) the label sequence of an infinite path of
\(X\). Then apply the label form, Theorem~\ref{thm:six-faces} or
Theorem~\ref{thm:determinization}.
\end{proof}

Part~(a) generalises two corpus theorems:
\begin{itemize}
\item \lean{DubC.digits\_eventually\_periodic\_of\_certOK},
\item \lean{DubC2.Gen.isEventuallyPeriodic\_of\_path}.
\end{itemize}
\cite[Lemma~2]{DN05} says that the carry word of a
nonzero \(\xi\) is never eventually periodic. So part~(c), applied to that word, is the shape of
every divisor-set theorem in the corpus: a certificate for the label form of \emph{some} cover
shows that no \(\xi\ne0\) avoids \(\mathcal P\).

\begin{theorem}[the obstruction principle]\label{thm:obstruction}\index{obstruction principle}
Let \(Y\) be finite and right-resolving, and suppose it realises forever through \(\sigma\), with
\(\sigma(u)\subseteq\bigcup_v\rho(v)\). If \(Y\) has a horseshoe at \(u\)
(Definition~\ref{def:horseshoe}) with two distinct petals, then no cover \(X\) of the same
system through \(\rho\) has the label form. If \(V\) is finite, no such \(X\) has zero entropy.
\end{theorem}

\begin{proof}
Theorem~\ref{thm:horseshoe} turns two petals into an infinite path \((s,\ell)\) of \(Y\) that
is not eventually periodic as a pair. Since \(Y\) is right-resolving, \(\ell\) is not eventually
periodic either (Proposition~\ref{prop:state-label} and its right-resolving converse).
Realisation gives a point of \(\sigma(u)\subseteq\bigcup\rho\) at which \(\ell\) is admissible
forever. Then apply Theorem~\ref{thm:transfer-per}\,(c).
\end{proof}

This is the principle behind the DubC2 horseshoe and branch witnesses, and behind
\cite[Thm.~42]{Oda26}. An obstruction found in one presentation, for instance a small
hand-made graph that realises its paths by the intermediate value theorem, rules out
certificates in \emph{every} presentation of the system.

\begin{theorem}[escape]\label{thm:transfer-escape}
Let \(V\) and \(A\) be finite and let \(X\) cover the system through \(\rho\).
\begin{enumerate}[label=(\alph*)]
\item If \(v\notin\mathrm{peel}_k(X)\) (Proposition~\ref{prop:peel}), no point of \(\rho(v)\)
  follows a word of length \(k\).
\item If the core of \(X\) is empty, no point of \(\bigcup_v\rho(v)\) follows a word of length
  \(|V|\). In particular no orbit from \(\bigcup_v\rho(v)\) is admissible forever.
\end{enumerate}
\end{theorem}

\begin{proof}
By Proposition~\ref{prop:shadow}\,(a), a followed word of length \(k\) is a walk of length
\(k\) from \(v\). Then apply Proposition~\ref{prop:peel} and Theorem~\ref{thm:core-empty}.
\end{proof}

This is the unavoidable-divisor-set theorem of \cite{DubOst06} and the escape bound of the
\texttt{Z32} funnels. The peel index of \(v\) bounds the escape time of every point of
\(\rho(v)\) (Theorem~\ref{thm:escape}).

\section{Counting itineraries}

The \emph{itineraries} of length \(n\) from a set \(T\subseteq S\) are the words followed by
some point of \(T\).

\begin{theorem}[entropy inequality]\label{thm:count}
Let \(V\) and \(A\) be finite and let \(X\) cover the system through \(\rho\).
\begin{enumerate}[label=(\alph*)]
\item The itineraries of length \(n\) from \(\rho(v)\) are words of \(X\) from \(v\), so there
  are at most \(\mathrm{npaths}(v,n)\) of them. If \(T\subseteq\bigcup_v\rho(v)\), there are at
  most \(\mathrm{ntot}(n)\) itineraries of length \(n\) from \(T\), and their exponential growth
  rate is at most \(h(X)\).
\item If \(X\) realises its words, the itineraries from \(\rho(v)\) are exactly the words from
  \(v\). If moreover \(X\) is right-resolving, then \(h(X)\) is the growth rate of the
  itineraries from the union of all \(\rho(v)\).
\end{enumerate}
\end{theorem}

\begin{proof}
(a) is Proposition~\ref{prop:shadow}\,(a) and Proposition~\ref{prop:count}. For (b), a
right-resolving graph has as many words as paths, so
\(\mathrm{ntot}(n)=\sum_v\#\mathrm{words}(v,n)\le|V|\cdot\#\{\text{itineraries}\}\), and a
constant factor does not change the growth rate.
\end{proof}

So \(h(X)\) bounds the complexity of the carry language of the system from above, through every
finite cover. An exact right-resolving cover, such as the \texttt{Z32} type certificate,
computes it.

\section{Dimension}

Let \(S\) be a (pseudo-)metric space and let \(K\subseteq S\) be a set of points of
\(\bigcup_v\rho(v)\) that follow words of every length. The itinerary \emph{localises} with
ratio \(r<1\) if two points of \(K\) in one \(\rho(v)\) that follow the same word of length
\(n\) are within \(c\,r^n\) of each other.

\begin{theorem}[dimension $\le$ entropy$/\log\lambda$]\label{thm:dimension}\index{dimension!Hausdorff, box}
Let \(V\) and \(A\) be finite, let \(X\) cover the system through \(\rho\), and let the itinerary
localise on \(K\) with ratio \(r=\lambda^{-1}\in(0,1)\). Then
\[
\overline{\dim}_B K\le\frac{h(X)}{\log\lambda},\qquad
\dim_H K\le\frac{h(X)}{\log\lambda}\quad(K\ne\emptyset).
\]
Localisation holds if every admissible step expands distances by at least \(\lambda\) and the run
endpoints of two such points are \(c\)-close.
\end{theorem}

\begin{proof}
Choose one point of \(K\) for each pair \((v,w)\) with \(w\) a word of length \(n\) from \(v\)
that some point of \(K\cap\rho(v)\) follows. The balls of radius \(c\,r^n\) around these points
cover \(K\). There are at most \(\sum_v\#\mathrm{words}(v,n)\le\mathrm{ntot}(n)\) of them. The
upper box dimension is the growth rate of the covering numbers at scales \(c\,r^n\), divided by
\(\log r^{-1}\) (\texttt{ForMathlib}'s \lean{upperBoxDim\_eq\_upperBoxDimWith} and
\lean{coveringGrowth\_const\_mul}). It is therefore at most \(h(X)/\log\lambda\). Finally,
\(\dim_H\le\overline{\dim}_B\). For the expanding criterion, induction along the run gives
\(d(x,y)\le r^{n}\,d(\mathrm{run}(x,w),\mathrm{run}(y,w))\).
\end{proof}

With the eigenvector bracket of Theorem~\ref{thm:brackets} this gives explicit numbers: a
super-eigenvector with ratio \(a/b\) yields \(\dim_H K\le\log(a/b)/\log\lambda\). That is the
inequality \lean{DubC2.Gen.upperBoxDim\_avoiders\_le} proves for one certificate format.
Chapter~\ref{ch:affine} derives it for every cover of the affine system (Theorem~\ref{thm:dim-avoiders}).

\section*{Notes}

Covers in this sense are one-sided, set-valued Markov partitions \cite[§6.5]{LM95}. A cover
need not be a partition, the sets \(\rho(v)\) may overlap, and \(X\) may have more arcs than the
system uses. The shadowing proposition replaces the coding map of a Markov partition by a choice
of path, and the transfer theorems need nothing more. The entropy inequality is the
graph-theoretic half of \cite[Prop.~4.1.13, ch.~4]{LM95}; the dimension bound is the counting
step of \cite[§3.2]{Fal03}.

\section*{Lean}

\begin{center}\small
\begin{tabular}{@{}>{\raggedright}p{0.3\textwidth}>{\raggedright\arraybackslash}p{0.64\textwidth}@{}}
\toprule
Text & Lean (\texttt{CarryGraph/Cover.lean}, \texttt{Dimension.lean})\\
\midrule
Definition~\ref{def:system} & \lean{run}, \lean{Follows}, \lean{orbit}, \lean{AdmForever},
  \lean{orbit\_succ'}, \lean{AdmForever.tail}, \lean{pre}, \lean{length\_pre},
  \lean{pre\_succ}, \lean{AdmForever.follows}\\
Definition~\ref{def:cover} & \lean{Cover}, \lean{Realizes}, \lean{RealizesInf}\\
Proposition~\ref{prop:shadow} & (a) \lean{Cover.follows}, \lean{Cover.isWord}; (b)
  \lean{Cover.exists\_iPath}\\
Constructions & \lean{Cover.mono}, \lean{Cover.mono\_adm}, \lean{Cover.map},
  \lean{Cover.comap}, \lean{Cover.labelProduct}, \lean{Cover.piProduct}, \lean{Cover.det},
  \lean{Cover.lineGraph}, \lean{Cover.ofSimulation}\\
Theorem~\ref{thm:transfer-per} & \lean{Cover.evPer\_of\_allPer},
  \lean{Cover.evPer\_of\_lnoBranch}, \lean{Cover.evPer\_of\_det}, \lean{Cover.not\_allPer}\\
Theorem~\ref{thm:obstruction} & \lean{Cover.not\_allPer\_of\_horseshoe},
  \lean{Cover.not\_lnoBranch\_of\_horseshoe}\\
Theorem~\ref{thm:transfer-escape} & \lean{Cover.not\_follows\_of\_not\_hasWalk},
  \lean{Cover.escape\_of\_peel}, \lean{Cover.escape\_of\_core\_empty},
  \lean{Cover.not\_admForever\_of\_core\_empty}\\
Theorem~\ref{thm:count} & (a) \lean{orbitWords}, \lean{orbitLang},
  \lean{Cover.orbitWords\_subset}, \lean{Cover.ncard\_orbitWords\_le},
  \lean{Cover.orbitLang\_subset}, \lean{Cover.ncard\_orbitLang\_le},
  \lean{Cover.orbitEntropy\_le}; (b) \lean{Realizes.words\_subset}, \lean{Cover.orbitWords\_eq},
  \lean{Cover.orbitEntropy\_eq}\\
Theorem~\ref{thm:dimension} & \lean{Cover.externalCoveringNumber\_le},
  \lean{Cover.upperBoxDim\_le}, \lean{Cover.dimH\_le}, \lean{edist\_le\_of\_follows},
  \lean{Cover.upperBoxDim\_le\_of\_expanding}\\
\bottomrule
\end{tabular}
\end{center}
